\chapter{Agent 协作流程与实现}

\section{公开角色与学生责任}

学生通过一个 Agent Runtime 使用五个角色。角色名称直接对应实验任务，模型 Provider 和具体 profile 由项目配置与任务类型决定。

\begin{table}[htbp]
  \centering
  \caption{学生可见 Agent 角色}
  \label{tab:agent-roles}
  \begin{tabularx}{\textwidth}{p{2.5cm}p{3.0cm}X}
    \toprule
    角色 & 项目权限 & 主要工作 \\
    \midrule
    ask & 只读问答 & 讨论体系结构、内核机制和课程资料，在使用知识库时给出来源 \\
    review & 只读评审 & 先执行 lint，再检查目标规格、相关规格和测试映射 \\
    implement & 临时工作树内写入 & 实现目标模块，生成测试并提交结构化结果 \\
    debug & 只读诊断 & 阅读最近运行记录，说明根因、证据和修复方向 \\
    verify & 只读检查 & 汇总公开测试、契约和 Spec ID 覆盖情况 \\
    \bottomrule
  \end{tabularx}
\end{table}

\code{agent config} 是模型设置命令。它记录 Provider、模型、base URL 和凭据环境变量名。聊天模型与知识库 embedding 分别配置，课程可以独立选择两类服务。当前实现支持 Anthropic、OpenAI、DeepSeek、Ollama 等主流服务提供商，也可以接入第三方 API。

\section{从问答到手写规格}

学生设计阶段采用一条清晰循环：

\begin{center}
\textbf{讨论} $\rightarrow$ \textbf{学生手写 Spec} $\rightarrow$ \textbf{spec lint} $\rightarrow$ \textbf{agent review} $\rightarrow$ \textbf{学生修改并提交}
\end{center}

\code{agent ask} 提供概念解释和方案比较。\code{spec lint} 用确定性程序检查 schema、引用、路径、等级和验证映射。\code{agent review} 读取 lint 结果、目标规格、相关模块与 \code{verifies} 关系，返回可以定位到具体字段的评审意见。三个命令都保留学生文件原状。学生根据讨论和评审亲手修改规格，并决定何时形成 Git 提交。

这条流程使 AI 参与设计讨论，同时让课程能够直接检查学生写下的系统方案、模块契约和接口选择。规格提交也成为后续实现的明确起点。

\section{Provider 与结构化结果}

Provider 层处理不同模型服务的请求协议、模型名称和凭据映射。角色 profile 定义任务提示、可用工具、读取范围和输出 schema。二者分离后，课程更换模型时仍然使用相同的 \code{owns}、测试字段和验证门禁。

需要机器处理的结果通过 MCP 提交工具交付。工具先验证 JSON schema，再检查路径、Spec ID、测试类型和任务状态。验证错误会直接返回同一模型线程，模型可以继续读取文件、修改实现并重新提交。VOS 只采用通过工具验证的结构化结果。

\section{代码生成与实现流程}

\code{vos agent implement <module>} 是学生使用的代码生成入口。任务开始时，平台读取已经提交的 ModuleSpec，将其中的操作、性质、不变量、接口依赖和测试要求整理成模型上下文，同时提供当前模块的源码、公共测试框架、已有测试目标和允许修改的文件范围。模型据此直接编写具体实现和配套测试，并提交可以执行的结果。

平台从目标模块的 \code{owns} 字段得到基本修改范围。确需调整其他模块时，学生要先提交 SpecPatch，说明修改原因、涉及模块和需要复查的性质，平台再把其中声明的文件加入本次任务。图 \ref{fig:agent-transaction} 展示代码生成、独立验证和写回学生项目的完整过程。

\begin{figure}[htbp]
  \centering
  \includegraphics[width=\textwidth]{agent-transaction.pdf}
  \caption{Agent implement 的临时工作树与验证流程}
  \label{fig:agent-transaction}
\end{figure}

实现结果的接受条件可以写成：

\begin{equation}
  \begin{aligned}
    accept ={}& clean_{base} \land diff \subseteq P_{allowed} \land targets_{valid} \\
              & {}\land build=pass \land checks_{module}=pass \land HEAD_{now}=HEAD_{base}
  \end{aligned}
\end{equation}

平台从当前提交创建临时 Git 工作区，模型只在其中编写代码和测试。学生正在使用的目录不会被未完成的修改覆盖。每轮生成结束后，VOS 从 Git 差异中读取实际修改的文件，再与允许范围比较。模型在回复中声明的修改列表只用于说明工作，不能代替这一步检查。

测试目标由模型通过结构化结果提出，VOS 负责检查并更新 \code{vos.yaml}。这样可以保证目标名称不与已有目标冲突，命令引用的测试文件真实存在，\code{verifies} 指向稳定的 Spec ID。随后，平台独立执行编译和全部已有、新增测试。代码通过检查后，VOS 在临时工作区创建 Git 提交，并在确认学生项目仍停留在任务开始时的提交后写回。提交记录本次运行编号和规格摘要，学生可以从代码提交查到生成时采用的规格与验证结果。

\section{自动生成测试与验证目标}

实现助手必须在生成内核代码的同时编写可执行测试。所谓自动生成测试，指模型根据 ModuleSpec 的操作、前后置条件、不变量和接口约束生成具体测试源码、驱动脚本与判定逻辑，再由平台实际运行。结构化结果至少包含 public、contract、fuzz、trace 四类测试目标和一组本地 hidden tests。

\begin{table}[htbp]
  \centering
  \caption{实现任务生成的测试材料}
  \begin{tabularx}{\textwidth}{p{3.0cm}X}
    \toprule
    类型 & 课程用途 \\
    \midrule
    public & 从用户或外部接口观察功能结果，例如系统调用返回值、串口输出和文件内容 \\
    contract & 检查 ModuleSpec、InterfaceSpec 中的前置条件、后置条件、错误行为和跨模块约定 \\
    fuzz & 用固定随机种子生成有限组输入，失败时保存触发输入和复现材料 \\
    trace & 运行限定规模的场景，通过明确的判定规则检查关键执行步骤或状态变化 \\
    hidden & 保存模型另行生成的本地补充用例，由 \code{verify --hidden} 显式执行 \\
    \bottomrule
  \end{tabularx}
\end{table}

每个测试目标都记录可执行程序、参数、工作目录、允许继承的环境变量、超时时间、对应的 Spec ID 和预期产物。平台在执行前检查路径是否位于项目内，测试文件是否存在，目标是否覆盖当前模块，随后才将其加入验证配置。

模糊测试的随机种子和用例数均写入目标。测试失败时，\code{reproduction\_artifact} 指向保存触发输入的文件，因此同一输入可以在调试阶段直接重放。固定种子和有限用例数也使课堂验收的运行时间与结果保持稳定。

实现阶段生成的 trace 测试是项目中长期保留的测试代码。它声明一个有界工作负载，例如限定次数的进程切换、页分配或文件操作，并给出判断结果是否正确的 oracle。oracle 可以检查最终状态、输出序列或关键事件的先后关系。测试每次都以同一工作负载运行，既能覆盖单次输入难以暴露的状态变化，也不会把无限运行或人工观察当作验收方式。

本地 hidden tests 的源码由结构化结果单独交付，不写入代码提交。平台记录这些用例所属的模块、生成它们的运行编号，以及规格、配置和测试内容的摘要。\code{vos verify --hidden} 只有在这些信息与当前项目一致时才执行用例，避免规格或实现改变后继续使用已经失效的测试。

\section{同一线程中的验证修正}

模型第一次提交的测试描述或实现可能因字段格式、文件路径、修改范围、编译错误或测试失败而被拒绝。运行时按照结构化结果、Git 差异、编译和测试四个层次生成简明诊断，再返回同一 Agent 线程，并保留当前临时工作区。模型能够读取实际错误，修改已有代码与测试，再次提交结果。平台始终重新执行检查，模型声称已经修复不能直接改变任务状态。通过后的提交来自这一连续修正过程，相关事件与命令输出一并进入本地运行记录。

\section{插桩调试与根因分析}

\code{vos agent debug} 从一次失败运行的编号或最近日志开始。平台先收集失败命令、退出状态、串口日志、构建产物和关联的 Spec ID，再让调试助手判断错误属于编译、启动、内存、并发、接口还是测试判定问题。输出采用固定结构，包含故障分类、可疑规格条款、证据链、下一步诊断命令和修复方向。

当已有日志不足以定位根因时，调试助手会生成一份临时插桩方案。方案包含 Git 补丁、测试刺激、普通输出的成功与失败模式，以及预期观察到的内部事件。平台只在单独创建的调试工作区中应用补丁，先检查补丁路径，再重新编译内核并执行每个诊断用例。插桩不会写回学生源码。

插桩代码用统一前缀 \code{VOS\_TRACE} 输出 JSON 事件。Runner 将串口日志中的普通程序输出与这些内部事件分离，分别保存完整串口记录、\code{trace.jsonl} 和单用例结果。一个诊断用例必须同时满足正常输出判定、没有命中失败模式、预期事件已经出现。仅打印一条预期事件不能被当成功能通过，这使内部执行路径和外部行为能够相互印证。

对于需要观察控制流和机器状态的问题，平台还为 QEMU 的 GDB 调试接口生成适配信息。调试助手可以设置断点，查看调用栈、寄存器、故障指令、源码位置和内存，并把 QEMU monitor 的只读状态作为补充证据。最终报告将日志、trace 事件、GDB 观察与相关规格放在同一证据链中，学生可以沿着``现象、内部路径、故障位置、修复方向''理解问题发生的原因。

调试插桩 trace 用于定位一次失败，任务结束后移除。实现阶段的 trace 测试随模块提交，并在后续 \code{vos verify} 中持续执行。

\section{确定性验证与 Agent 复核}

平台将实际验收与模型解释分成两个步骤。\code{vos verify} 由确定性 Runner 执行，要求项目处于已提交状态，依次检查规格、编译内核，并运行 \code{vos.yaml} 中的 public、contract、固定种子 fuzz 和有界 trace 目标。每个目标的退出状态、超时、日志和产物共同决定结果，整个过程不调用模型。

\code{vos agent verify} 先在临时工作区执行同一套公开验证，再把实际结果、五类规格、测试到 Spec ID 的映射和现有证据交给只读验证助手。助手逐项寻找测试结果之间的矛盾，检查每个稳定 Spec ID 是否具有有意义的 public、contract、fuzz 或 trace 覆盖，并给出证据引用和补充建议。结构化输出中的验证状态必须与 Runner 的结果完全一致，Agent 不能把失败改写为通过，也不能推测本地 hidden tests 的内容。

平台在 \code{agent debug}、\code{agent verify}、\code{agent review} 和 \code{agent ask} 运行前后比较 Git 文件指纹。出现源码或规格变化时，任务会被判为违反只读约定。由此，调试和复核负责解释证据，代码与测试的改变仍然回到明确的实现流程中完成。
